\chapter{Introduction}

\section{Background}

This chapter shows the basic document structure. A superscript numerical citation is produced by \verb|\cite|, for example \cite{sample-journal}. An inline citation is produced by \verb|\inlinecite|, for example \inlinecite{sample-book}.

The first paragraph after a heading is indented. Equations are numbered by chapter, as shown in equation \eqref{eq:model-example}.
\begin{equation}\label{eq:model-example}
  \min_{x\in\mathbb{R}^{d}} F(x) := f(x)+g(Bx).
\end{equation}

\section{Figures and Tables}

Figure \ref{fig:placeholder} shows a placeholder figure. Replace the framed box with \verb|\includegraphics| when a real figure file is available.
\begin{figure}[htbp]
  \centering
  \fbox{\rule{0pt}{35mm}\rule{80mm}{0pt}}
  \caption{A placeholder figure}
  \label{fig:placeholder}
\end{figure}

Table \ref{tab:notation} shows a standard three-line table.
\begin{table}[htbp]
  \centering
  \caption{Basic notation}
  \label{tab:notation}
  \begin{tabular}{ll}
    \toprule
    Symbol & Meaning \\
    \midrule
    $x$ & Primal variable \\
    $y$ & Dual variable \\
    $B$ & Linear operator \\
    \bottomrule
  \end{tabular}
\end{table}

\section{Section Levels}

\subsection{Subsection Example}

This is a subsection. Its numbering follows the form 1.3.1.

\subsubsection{Subsubsection Example}

This is a subsubsection. Its numbering follows the form 1.3.1.1.
